%============================================================
% Part 5 opener  (generated by tools/outline.py; design: preamble/partpage.tex)
%============================================================
\dppartopener{Bits in Problem Solving}{CF 1500 $\to$ CF 2600}{warningred}%
  {Randomised hashing, counting over binary prefixes, carries, interactive and constructive problems built on bits.}

\dppartstats{5}{20}{40}{150}

\begin{dpjourney}
\dpstep{23}{Randomisation and Hashing with Bits}{Robust hashing $\cdot$ Random masks $\cdot$ Hashing states $\cdot$ Probabilistic structures}{CF 1700 -- 2600}{no}
\dpstep{24}{Counting over Binary Prefixes}{Free suffixes $\cdot$ Bounded XOR $\cdot$ Sums over ranges $\cdot$ Counting pairs under a bound}{CF 1600 -- 2500}{no}
\dpstep{25}{Arithmetic Meets Bitwise: Carries and Borrows}{No-carry additions $\cdot$ Borrows $\cdot$ Mixed inequalities $\cdot$ Carry chains}{CF 1700 -- 2600}{no}
\dpstep{26}{Interactive Bit Problems}{Sums from OR and AND $\cdot$ XOR queries $\cdot$ Bit by bit $\cdot$ Budgets and adversaries}{CF 1700 -- 2600}{no}
\dpstep{27}{Constructive Bit Problems}{Fixing the total XOR $\cdot$ High-bit reservations $\cdot$ Bitwise permutations $\cdot$ Impossibility proofs}{CF 1500 -- 2600}{no}
\end{dpjourney}

\dppartgoals{
  \item hash states and multisets with random 64-bit values and robust mixers;
  \item count numbers up to $n$ with a bitwise property by fixing a binary prefix;
  \item reason about carries when arithmetic and bitwise operations meet;
  \item recover hidden values with bitwise queries, and construct arrays with XOR, AND and OR constraints.
}
